% !TEX program = xelatex
\documentclass[12pt,a4paper]{ctexart}

\usepackage{geometry}
\geometry{left=2.5cm,right=2.5cm,top=2.5cm,bottom=2.5cm}
\usepackage{tikz}
\usetikzlibrary{shapes.geometric, arrows.meta, positioning, fit, backgrounds}
\usepackage{amsmath}
\usepackage{hyperref}
\usepackage{fontawesome5}
\usepackage{listings}
\usepackage{xcolor}

\definecolor{boxblue}{HTML}{4472C4}
\definecolor{boxorange}{HTML}{ED7D31}
\definecolor{boxgreen}{HTML}{70AD47}
\definecolor{boxred}{HTML}{C0504D}
\definecolor{boxgray}{HTML}{808080}
\definecolor{boxpurple}{HTML}{9B59B6}
\definecolor{bg}{HTML}{F2F2F2}

\tikzstyle{module} = [rectangle, draw, rounded corners, minimum width=2.5cm, minimum height=1cm, text centered, font=\small, align=center]
\tikzstyle{data} = [rectangle, draw, fill=bg, minimum width=2cm, minimum height=0.7cm, text centered, font=\footnotesize, align=center]
\tikzstyle{gate} = [diamond, draw, fill=boxred!10, aspect=2, minimum width=1.8cm, minimum height=0.8cm, text centered, font=\footnotesize, align=center]
\tikzstyle{arrow} = [thick, ->, >=stealth]

\title{\textbf{供电设备图像异动检测系统}\\[4pt]
       \large 流程说明与模块拆解}
\author{}
\date{\today}

\begin{document}
\maketitle
\tableofcontents
\newpage

%===========================================================================
\section{系统总览：一条图的两头}
\label{sec:overview}

系统的任务很简单：业务方拍了两张图——一张叫\textbf{标准源图}（这个设备正常时长什么样），一张叫\textbf{实时巡检图}（现在长什么样），另外在标准图上用 YOLO 框标出了要检查的部件。系统要做的事只有三件：

\begin{enumerate}
\item 先把两张图对上——实时图里的那个位置，到底对应标准图的哪个位置？
\item 对上了之后，把两个状态的同一片区域拿出来比——有没有哪里不一样？
\item 给出结论：没变化（高置信）、有变化（框出来）、还是说"我判断不了"（告知原因）。
\end{enumerate}

整个流程是一条直线，分成\textbf{七个模块}，每个模块只做一件事，做完把结果交给下一个模块。如下图所示。

\begin{figure}[htbp]
\centering
\resizebox{\textwidth}{!}{%
\begin{tikzpicture}[node distance=0.5cm]
\node (in1) [data] {标准源图.jpg};
\node (in2) [data, below=0.2cm of in1] {对比截图.jpg};
\node (in3) [data, below=0.2cm of in2] {标准源图坐标.txt (YOLO)};
\node (bg1) [draw, fit=(in1)(in2)(in3), fill=bg, rounded corners] {};

\node (m1) [module, fill=boxblue!20, right=1.5cm of bg1] {1. 数据准入};
\node (m2) [module, fill=boxorange!20, right=1.5cm of m1] {2. 配准级联};
\node (m3) [module, fill=boxpurple!20, right=1.5cm of m2] {3. 部件映射};
\node (gate1) [gate, right=1.5cm of m3] {通过？};
\node (m4) [module, fill=boxred!20, right=1.2cm of gate1] {4. 差异检测};
\node (m5) [module, fill=boxgreen!20, right=1.5cm of m4] {5. 证据产物};
\node (m6) [module, fill=boxgray!20, right=1.5cm of m5] {6. 输出结论};

\node (out1) [data, right=1.5cm of m6] {对齐图 / mask};
\node (out2) [data, below=0.2cm of out1] {差异图 / 热力图};
\node (out3) [data, below=0.2cm of out2] {标注图 / 结构化结果};

\node (unavail) [module, fill=boxred!10, below=1.5cm of gate1] {检测不可用\\+ 证据产物};

\draw [arrow] (bg1) -- (m1);
\draw [arrow] (m1) -- (m2);
\draw [arrow] (m2) -- (m3);
\draw [arrow] (m3) -- (gate1);
\draw [arrow] (gate1) -- node[above] {\footnotesize 是} (m4);
\draw [arrow] (gate1) -- node[left] {\footnotesize 否} (unavail);
\draw [arrow] (m4) -- (m5);
\draw [arrow] (m5) -- (m6);
\draw [arrow] (m6) -- (m6 -| out1);
\end{tikzpicture}%
}
\caption{系统主流程（单部件）}
\label{fig:pipeline}
\end{figure}

六个模块对应六个源文件，再加一个负责"多个部件凑一张图"的聚合层（模块7），构成了整个系统。下面按顺序逐个拆解。


%===========================================================================
\section{模块1：数据准入 —— 把文件读进来、数读对}

\subsection*{对应文件}
\begin{itemize}
\item \texttt{references.py} —— 案例目录发现、图片读取
\item \texttt{roi.py} —— YOLO坐标文件解析
\end{itemize}

\subsection*{做了什么事}

\begin{enumerate}
\item 扫描一个案例目录，找到三样东西：
  \begin{itemize}
  \item \texttt{对比截图.jpg} —— 实时巡检拍的那张
  \item \texttt{标准源图.jpg} —— 这个机位的标准参照
  \item \texttt{标准源图坐标.txt} —— YOLO框文件，每行一个部件：\\[2pt]
    \texttt{类别ID\qquad 中心x\qquad 中心y\qquad 宽度\qquad 高度}\\[2pt]
    坐标都归一化到$[0,1]$（即0.5表示画面正中间，1.0表示最右/最下）
  \item 可选：\texttt{新增标准源图0.jpg}, \texttt{新增标准源图1.jpg}, ... —— 其他时刻拍的标准图（比如晚上拍的）
  \end{itemize}
\item 图片读取走的是 \texttt{numpy.fromfile + cv2.imdecode} 路径，而不是直接用 OpenCV 的 \texttt{imread}——因为 Windows 中文路径 OpenCV 不认
\item YOLO坐标文件解析时做了容错处理：坐标从$[-0.01, 1.01]$都在容忍范围内（稍微出框一点可以裁回来），超出这个范围才算格式非法
\item 解析出来的每行坐标转成像素矩形：$x = (cx - w/2) \times W$, $y = (cy - h/2)\times H$, $w = r_w \times W$, $h = r_h \times H$（其中 $W, H$ 是标准图的宽高，$r_w, r_h$ 是归一化的宽高）
\end{enumerate}

\subsection*{输入 → 输出}

\noindent\begin{tabular}{|p{4cm}|p{4cm}|p{5cm}|}
\hline
\textbf{输入} & \textbf{输出} & \textbf{说明} \\
\hline
案例目录路径 & \texttt{CaseInput} 对象 & 包含实时图路径、标准图路径列表、YOLO文件路径 \\
\hline
YOLO坐标文件 & \texttt{list[Roi]} & 每行解析成一个Roi对象（类别、中心x/y、宽、高，全归一化） \\
\hline
图片路径 & \texttt{np.ndarray} (BGR) & 三维数组，H×W×3，uint8 \\
\hline
归一化Roi + 图尺寸 & \texttt{(x, y, w, h)} & 像素坐标的矩形，整数 \\
\hline
\end{tabular}


%===========================================================================
\section{模块2：配准级联 —— 把两张图对上}

\subsection*{对应文件}
\begin{itemize}
\item \texttt{alignment.py} —— 完整的对齐流水线+三态诊断分类
\item \texttt{config.py} 中 \texttt{AlignmentThresholds} 部分 —— 所有门槛都来自配置，不写死
\end{itemize}

\subsection*{做了什么事}

配准级联是一条"先快后稳"的串联流水线。它先用ORB完成一次完整检查；只有ORB给出不可靠或不可用、且配置明确允许时，才让SuperPoint+LightGlue补一次全图匹配点。补点不是放行：补出的点仍要重新走RANSAC、投影几何、有效重叠和ECC微调这几道原门槛。默认配置关闭这条补路，待依赖许可审批完成后才可在日常检测中打开。

\textbf{第1步：ORB特征提取}\\[2pt]
两张灰度图各提2000个ORB特征点（Oriented FAST + Rotated BRIEF，一种经典的、不需要GPU的特征检测算法）。每个点附着一个256位的二进制描述符。

\textbf{第2步：特征匹配}\\[2pt]
用暴力匹配器（BFMatcher）把标准图的描述符跟实时图的比，选"最近邻距离除以次近邻距离 $< 0.75$"的点对留下来——比值越小，匹配越可靠。

\textbf{第3步：RANSAC求单应矩阵}\\[2pt]
拿匹配点的坐标对，用RANSAC（随机抽样一致）算法估计一个$3\times 3$的单应矩阵$H$。它的作用是把标准图上的坐标映射到实时图上：$[x_{\text{live}}, y_{\text{live}}, 1]^T \sim H \cdot [x_{\text{std}}, y_{\text{std}}, 1]^T$。RANSAC会同时把"跟大多数点不一致的"点标为外点，只留内点。

这一步算出来的几个关键数字：内点数量、内点率（内点/匹配点）、重投影误差（内点被$H$映射后的平均像素偏差）、空间覆盖度（内点占标准图面积的比例）。

\textbf{第4步：算投影几何}\\[2pt]
把标准图的四个角通过$H$投到实时图里，算：中心位移（两个图的画面中心差了多少像素）、角点位移（四个角各差多少）、投影面积比（投过去的面积占实时图的比例）、有几个角落在实时图画面内。

\textbf{第5步：Warp + 有效mask}\\[2pt]
把实时图按照$H^{-1}$（反变换）扭曲到标准图的坐标系里：标准图每个像素去找实时图上对应的位置，用双线性插值填颜色。扭曲之后会有黑边（实时图够不到的位置），这些地方在"有效mask"里标为无效。

\textbf{第6步：ECC微调}\\[2pt]
warp之后，两图可能还有几个像素的对不齐。用OpenCV的ECC（增强相关系数）算法在局部做一个仿射微调——在有效mask范围内，找一个能让两图灰度值更一致的小幅平移+旋转+缩放。收敛了就得到更高的对齐精度；如果50次迭代还不收敛，记录"ECC未收敛"并把这个事实当作负证据传给后面的门禁。

\textbf{第7步：近恒等+父子图检测}\\[2pt]
如果$H$近乎恒等变换（中心位移不到画面对角线的5\%，面积比在$[0.85,1.15]$之内）——说明两张图的拍摄位置基本没动。\textbf{但是}，如果同一个$H$下内点覆盖度低于0.2且ECC还没收敛——这就是典型的"父子图"场景：标准图是局部放大，实时图是广角，虽然$H\approx I$看起来不错（因为特征点凑巧在地面纹理上匹配上了），但实际上两图的尺度完全对不上。这种情况要专门打上标记（\texttt{suspected\_zoom\_parent\_child\_mismatch}），后面的门禁会据此拦下。

\textbf{第8步：三态诊断分类}\\[2pt]
综合前面1--7步的所有指标，给出三态结论：
\begin{itemize}
\item \textbf{usable（可用）}：所有指标都在配置的门槛之上，可以继续往下走
\item \textbf{unreliable（不可靠）}：虽然还能算出一个变换，但有关键证据不达标。日常检测到这里就停止，统一输出"检测不可用"；后面即使为了排查而算出了部件指标，也不能把它变成"无异动"、"部件差异"或"图像不匹配"
\item \textbf{unavailable（不可用）}：根本配不出变换，或者投影几何非法——直接终止，报"检测不可用"
\end{itemize}

\textbf{ORB失败时的辅助尝试}\\[2pt]
SuperPoint是一种从图像中找稳定特征点的模型，LightGlue是一种把两张图的特征点配成对的模型。它们只回答"这两个像素位置可能对应"，不回答"两张图已经可信对齐"，更不会直接给部件画框。模型输出的匹配点会像ORB点一样重新交给第3--8步处理；任何一项门槛不过，结果仍是"检测不可用"。同一对标准图和实时图在一个多部件案例中只计算一次对齐结果，后续部件复用该结果，避免重复推理。

\begin{figure}[htbp]
\centering
\begin{tikzpicture}[node distance=0.6cm]
\node (s1) [module, fill=boxblue!10, minimum width=3cm] {ORB特征提取};
\node (s2) [module, fill=boxblue!10, minimum width=3cm, below=of s1] {BFMatcher匹配\\+ 比例测试};
\node (s3) [module, fill=boxorange!10, minimum width=3cm, below=of s2] {RANSAC单应矩阵$H$\\+ 内点/误差/覆盖度};
\node (s4) [module, fill=boxorange!10, minimum width=3cm, below=of s3] {投影几何\\中心/角点位移、面积比};
\node (s5) [module, fill=boxpurple!10, minimum width=3cm, below=of s4] {Warp + 有效mask};
\node (s6) [module, fill=boxpurple!10, minimum width=3cm, below=of s5] {ECC仿射微调};
\node (s7) [module, fill=boxred!10, minimum width=3cm, below=of s6] {近恒等+父子图检测};
\node (s8) [module, fill=boxgreen!10, minimum width=3cm, below=of s7] {三态诊断\\usable/unreliable/unavailable};
\draw [arrow] (s1) -- (s2); \draw [arrow] (s2) -- (s3);
\draw [arrow] (s3) -- (s4); \draw [arrow] (s4) -- (s5);
\draw [arrow] (s5) -- (s6); \draw [arrow] (s6) -- (s7); \draw [arrow] (s7) -- (s8);
\end{tikzpicture}
\caption{配准级联的主流程；ORB不可靠时，可选的SuperPoint+LightGlue只补匹配点，随后仍回到第3--8步}
\label{fig:alignment}
\end{figure}

\subsection*{输入 → 输出}
\noindent\begin{tabular}{|p{4cm}|p{4cm}|p{5cm}|}
\hline
\textbf{输入} & \textbf{输出} & \textbf{说明} \\
\hline
标准灰度图、实时灰度图、\texttt{CalibratedConfig} & \texttt{AlignmentResult} & 含$H$矩阵、$H^{-1}$、aligned\_live图、valid\_mask、全部配准指标、诊断分类 \\
\hline
\end{tabular}


%===========================================================================
\section{模块3：部件映射门禁 —— 对上之后，再确认"真能用"}

\subsection*{对应文件}
\begin{itemize}
\item \texttt{mapping.py} —— 四道组件级检查
\end{itemize}

\subsection*{做了什么事}

配准给出的是一个"全图"级别的对齐关系。但系统要比较的是\textbf{某一个具体部件}的ROI。所以需要再走一遍跟部件直接相关的检查——看看这个部件在实时图里到底能不能拿出来比。

检查顺序是固定的，前一道不过就不走后一道（短路）：

\textbf{第1道：候选定位（B层）}\\[2pt]
把标准图上的ROI（要比较的那一小块区域）四个角通过$H$投射到实时图里，看投过去的四边形有多少在画面内。当前正式配置要求至少70\%在画面内。达不到说明这个位置已经不能可靠对齐，报"配准失败"。

\textbf{第2道：有效像素重叠（B层）}\\[2pt]
在ROI区域内，有效mask覆盖了多少像素？当前正式配置要求至少85\%。达不到说明要么是扭曲后的黑边占了太大面积，要么是配准不稳——报"配准失败"。

\textbf{第3道：有效分辨率（C层）}\\[2pt]
根据$H$在部件附近的缩放关系，算出：实时图上这个部件到底占了多少个像素？当前正式配置要求有效采样尺度至少0.4倍（\texttt{effective\_resolution\_scale\_min}），有效宽和高都至少48像素。达不到说明实时图里的部件太小或太糊，放大图片也不会凭空长出细节，报"匹配不确定"。

\textbf{第4道：配准后表观一致性（D层）}\\[2pt]
这是最关键的一道。在ROI里用有效mask内那些"配得上"的像素，算两个数：

\begin{itemize}
\item \textbf{NCC（归一化互相关）}：\[ \text{NCC} = \frac{\sum(A_i - \bar{A})(B_i - \bar{B})}{\sqrt{\sum(A_i-\bar{A})^2} \cdot \sqrt{\sum(B_i-\bar{B})^2}} \] 范围$[-1,1]$。1表示两张图的灰度"变化走势"完全一致；0表示互不相关；负数表示方向反了（标准图暗的地方实时图反而亮）
\item \textbf{SSIM（结构相似性）}：在每个11×11的小窗口里同时看平均亮度、明暗反差和纹理是否同向，再把所有窗口的结果取平均：\[\operatorname{SSIM}(A,B)=\frac{(2\mu_A\mu_B+C_1)(2\sigma_{AB}+C_2)}{(\mu_A^2+\mu_B^2+C_1)(\sigma_A^2+\sigma_B^2+C_2)}\] 其中$\mu$是平均亮度，$\sigma$是明暗变化程度，$\sigma_{AB}$表示两张图是否一起变亮或变暗；$C_1,C_2$是防止分母为零的小常数。范围大致在$[-1,1]$，越接近1越像。
\end{itemize}

原始NCC不到0.35，或原始SSIM不到0.2，只说明"这两块看起来不像"，可能是机位拍错，也可能只是日夜、曝光或阴影不同。这时先不进入差异检测，而是做三次复核：把两图的亮度分布拉到相近后再算NCC；去掉局部明暗起伏后再算NCC；只看边缘强弱后再算NCC。

\textbf{图像不匹配的确认规则}：全图配准和前三道部件检查必须全部通过；原始NCC也必须低于0.35；并且上面三次复核的NCC都仍低于0.35。只有同时满足这些条件，才输出"图像不匹配异动"。如果任一次复核还能对上，或者只有SSIM低，则证据还不够，输出"检测不可用"，绝不往下走比较。

\textbf{还有一个额外逻辑}：映射门禁同时算了一个\texttt{comparison\_scale}（比较用的缩放倍率），取 $\min(1.0, S)$（$S$是有效分辨率尺度）——如果实时图比标准图糙（分辨率低），就把两图都缩小到实时图的分辨率再比，不靠上采样造假。

\subsection*{输入 → 输出}
\noindent\begin{tabular}{|p{4cm}|p{4cm}|p{5cm}|}
\hline
\textbf{输入} & \textbf{输出} & \textbf{说明} \\
\hline
标准BGR图、warp后的实时图、有效mask、$H$矩阵、ROI矩形、实时图尺寸、配置 & \texttt{ComponentMappingEvidence} & 含四道结果、failure\_reason、comparison\_scale \\
\hline
\end{tabular}


%===========================================================================
\section{模块4：差异检测 —— 真·找不同}

\subsection*{对应文件}
\begin{itemize}
\item \texttt{compare.py} —— 抗光照颜色差、局部结构差、图像清晰度检查和候选分级
\end{itemize}

\subsection*{做了什么事}

模块3全部通过之后，两图已经对齐好了，ROI内的像素才允许进入"找不同"。这里不再直接拿原图逐像素相减，而是分别看颜色和结构，并先判断实时图是不是已经糊到不能可靠比较。

\textbf{通道1：先压掉光照影响，再看颜色}\\[2pt]
BGR先转Lab颜色空间：$L$表示明暗，$a$、$b$表示两条颜色轴。先在有效ROI内把实时图的$L$亮度分布拉到与标准图相近，再计算颜色距离：\[ \Delta E = \sqrt{(L_1-L_2)^2 + (a_1-a_2)^2 + (b_1-b_2)^2} \]
这里的$\Delta E$就是两种颜色在Lab空间里的直线距离。全图亮度中位数差得越多，颜色通道的权重越低，但不会低于配置下限。这样日夜、曝光和阴影不会轻易变成颜色异动；局部的油污、变色或涂层变化仍会留下信号。

\textbf{通道2：看局部结构是否变了}\\[2pt]
两图都先减去周围一大片区域的平均亮度，去掉缓慢变化的阴影；再用Sobel边缘算子分别求水平和竖直方向的边缘。Sobel不是缩写，它是一种用周围3×3像素估计边缘方向的算法。结构分数主要看两张图边缘方向是否一致：\[ s_{\text{structure}}=\frac{1-\cos\theta}{2}\times \min\left(1,\frac{\max(G_1,G_2)}{T_G}\right) \]
其中$\theta$是两张图在同一像素处的边缘方向夹角，$G_1,G_2$是边缘强度，$T_G$是配置中的强边缘尺度。方向一致时分数接近0；一个边缘消失、新增或方向明显改变时分数升高。这样比单纯比较边缘强弱更不容易把曝光变化误当成裂纹或缺件。

\textbf{先检查清晰度}\\[2pt]
在出业务结论前，系统用Laplacian（二阶边缘变化）计算标准图和实时图在有效ROI内的清晰度，再算较小值与较大值的比：\[ q=\frac{\min(S_{\text{std}},S_{\text{live}})}{\max(S_{\text{std}},S_{\text{live}})} \]
当$q$低于配置下限时，说明实时图相对标准图明显变糊。此时仍保留差异候选和五张证据图，但日常检测统一返回"检测不可用"，不会把拍糊了误报成设备异动，也不会假装"无异动"。

\textbf{合并与候选生成}\\[2pt]
两条通道的分数图取逐像素最大值，得到一张"融合差异热力图"。在这张热力图上：
\begin{enumerate}
\item 大于判定门槛（开发值=0.5）的位置标为1，其余标为0
\item 形态学操作（开运算）去噪
\item 找轮廓，把每个连通区域变成一个候选框$(x, y, w, h)$
\item 算每个候选的面积和置信度。置信度由区域内热力图平均值和面积共同得到；小区域不会因一个亮点就拿到很高把握
\item 所有候选都保留为原始证据。面积过小或置信度过低的候选会标成"仅供复核"，画成橙色框；只有同时达到面积和置信度门槛的候选才允许触发业务异动，画成红色框
\end{enumerate}

\subsection*{输入 → 输出}
\noindent\begin{tabular}{|p{4cm}|p{4cm}|p{5cm}|}
\hline
\textbf{输入} & \textbf{输出} & \textbf{说明} \\
\hline
标准BGR、warp后实时BGR、有效mask、ROI矩形、配置、比较倍率 & \texttt{CompareResult} & 含二值差异mask、float32热力图、清晰度证据、全部候选和可触发结论的候选 \\
\hline
\end{tabular}


%===========================================================================
\section{模块5：证据产物 —— 每一次判定都要能翻出来查}

\subsection*{对应文件}
\begin{itemize}
\item \texttt{artifacts.py} —— 五类证据图的生成
\end{itemize}

\subsection*{做了什么事}

不管最终结论是"无异动""有异动"还是"检测不可用"，每次都产出同样的五张图，存在调用方指定的目录里。这样以后有人想质疑"为什么判成这样"，可以直接翻证据。

\begin{enumerate}
\item \textbf{alignment.png} （对齐证据图）—— 左半边是标准图，右半边是实时图，中间用绿色线连的是标准图四角投影到实时图的位置，黄色线连的是RANSAC内点对（最多120根，间隔采样保证看得清）
\item \textbf{valid\_mask.png} （有效区域图）—— 白色表示这块像素可以比，黑色表示warp黑边或者无效区。二值图
\item \textbf{difference\_mask.png} （差异区域图）—— 白色表示被判定为"有差异"的像素。在整图大小上画，方便跟原图对着看
\item \textbf{difference\_heatmap.png} （差异热力图）—— 彩色热力图，越红表示差异分数越高。也是整图大小
\item \textbf{annotated.png} （标注图）—— 在实时图上画框：绿色框是ROI投影过去的位置；红色框是可触发结论的候选，橙色框是仅供复核的原始候选；候选框旁写该区域的置信度。图中没有颜色条
\end{enumerate}

如果流程在中途就被拦下（比如配安装不上、映射没过），照样生成对齐图和有效区域图——只是差异图是空的。这样"为什么没比出来"也是可以验证的，不会出现"没结果"但"不知道是不是根本没跑"的情况。

\subsection*{输入 → 输出}
\noindent\begin{tabular}{|p{4cm}|p{4cm}|p{5cm}|}
\hline
\textbf{输入} & \textbf{输出} & \textbf{说明} \\
\hline
标准图、实时图、$H$矩阵、配准结果、差异结果、候选框、输出目录 & 5个PNG文件路径 & \texttt{ArtifactSet}，每个都是本地路径 \\
\hline
\end{tabular}


%===========================================================================
\section{模块6：服务层 —— 把1--5串起来}

\subsection*{对应文件}
\begin{itemize}
\item \texttt{service.py} —— \texttt{ImageComparisonService}，单部件对比的唯一入口
\end{itemize}

\subsection*{做了什么事}

service就是"调度员"。它负责把模块1--5按正确的顺序和分支逻辑串起来：

\begin{enumerate}
\item 读图、读ROI（调模块1）
\item 跑配准（调模块2）
\item 如果全图配准诊断不是usable $\to$ 写不可用证据，返回"检测不可用"。为了排查问题可以保留已算出的部件指标，但它们不能进入业务判断
\item 全图配准可用时，才把实时图扭到标准图坐标，并跑部件映射（调模块3）
\item 映射确认"图像不匹配" $\to$ 写证据，返回"图像不匹配异动"
\item 其他映射不通过 $\to$ 写不可用证据，返回"检测不可用"（带原因）
\item 跑差异检测（调模块4）
\item 把候选区域从标准图坐标系反投回实时图（用$H$）
\item 写证据产物（调模块5）
\item 若清晰度不足 $\to$ "检测不可用"；否则有可触发结论的候选 $\to$ "检出异动"，只有没有这类候选时才是"高置信无异动"。原始小候选仍会保留在结果和证据图里，但不会单独触发异动
\end{enumerate}

\textbf{两种调用方式}：
\begin{itemize}
\item \texttt{compare()} —— 日常检测模式：必须传版本化的标定配置，返回带业务结论的 ComparisonResult
\item \texttt{observe()} —— 标定观察模式：可用开发配置，返回 CalibrationComponentObservation，\textbf{不含业务结论}（没有 normal/anomaly 字段）
\end{itemize}

\subsection*{输入 → 输出}
\noindent\begin{tabular}{|p{4cm}|p{4cm}|p{5cm}|}
\hline
\textbf{输入} & \textbf{输出} & \textbf{说明} \\
\hline
标准图路径、实时图路径、YOLO字符串(\texttt{"17 0.5 0.5 0.4 0.5"})、输出目录、配置路径 & \texttt{ComparisonResult}（日常）/ \texttt{CalibrationComponentObservation}（标定） & 日常含四类结论；标定只含观察记录和全部证据产物路径 \\
\hline
\end{tabular}


%===========================================================================
\section{模块7：多部件聚合 —— 一张图多个部件怎么办}

\subsection*{对应文件}
\begin{itemize}
\item \texttt{multi\_component.py} —— \texttt{MultiComponentImageComparisonService}
\end{itemize}

\subsection*{做了什么事}

一张标准图上的YOLO文件里有多个部件，系统要先选一个"最好对齐"的参考图（可能有多张标准图），然后对每个部件挨个跑模块6。

\textbf{选参考图的逻辑}：
\begin{enumerate}
\item 对每张参考图（主标准图 + 可选的新增标准图0/1/2...），跟实时图跑一次配准
\item 给每张参考图打分：usable优先 > unreliable > unavailable；同档内比内点率、重投影误差、有效重叠
\item 选分数最高的那张作为"选中参考图"
\item \textbf{所有部件都用同一张参考图}——因为后续的整图结论需要一个统一的坐标系
\end{enumerate}

\textbf{聚合结论的逻辑}：
\begin{itemize}
\item 任一部件"图像不匹配异动" $\to$ 整图"图像不匹配异动"
\item 否则，任一部件"部件差异异动" $\to$ 整图"部件差异异动"
\item 没有异动、但任一部件"检测不可用" $\to$ 整图"检测不可用"
\item 所有部件全部"高置信无异动" $\to$ 整图"高置信无异动"
\end{itemize}

\subsection*{输入 → 输出}
\noindent\begin{tabular}{|p{4cm}|p{4cm}|p{5cm}|}
\hline
\textbf{输入} & \textbf{输出} & \textbf{说明} \\
\hline
案例目录 & \texttt{CalibrationObservation}（标定）/ \texttt{ImageComparisonResult}（日常） & 含选中的参考图、每个部件的结论/观察、整图聚合结论、manifest JSON \\
\hline
\end{tabular}


%===========================================================================
\section{双模式流程图对比}

\begin{figure}[htbp]
\centering
\begin{tikzpicture}[node distance=1.2cm]
% Left: calibration
\node (ctitle) [font=\bfseries] at (0,0) {标定模式（Calibration）};
\node (c1) [module, fill=boxblue!20, below=0.5cm of ctitle, minimum width=3.5cm] {读入案例 (references.py + roi.py)};
\node (c2) [module, fill=boxorange!20, below=of c1, minimum width=3.5cm] {每张参考图跑配准 (alignment.py)};
\node (c3) [module, fill=boxorange!20, below=of c2, minimum width=3.5cm] {选最佳参考图};
\node (c4) [module, fill=boxpurple!20, below=of c3, minimum width=3.5cm] {逐部件跑 observe()};
\node (c5) [module, fill=boxgreen!20, below=of c4, minimum width=3.5cm] {产出标定观察 JSON + 证据图};
\node (c6) [data, fill=bg, below=of c5, minimum width=3.5cm] {无业务结论\\只产指标和证据};

\draw [arrow] (c1) -- (c2); \draw [arrow] (c2) -- (c3);
\draw [arrow] (c3) -- (c4); \draw [arrow] (c4) -- (c5); \draw [arrow] (c5) -- (c6);

% Right: daily
\node (dtitle) [font=\bfseries] at (7.5,0) {日常检测模式（Daily Detection）};
\node (d0) [data, fill=boxred!10, below=0.5cm of dtitle, minimum width=3.5cm] {\textbf{先加载标定配置}\\\textbf{拒绝开发配置}};
\node (d1) [module, fill=boxblue!20, below=of d0, minimum width=3.5cm] {读入案例  (references.py + roi.py)};
\node (d2) [module, fill=boxorange!20, below=of d1, minimum width=3.5cm] {每张参考图跑配准 (alignment.py)};
\node (d3) [module, fill=boxorange!20, below=of d2, minimum width=3.5cm] {选最佳参考图};
\node (d4) [module, fill=boxpurple!20, below=of d3, minimum width=3.5cm] {逐部件跑 compare()};
\node (d5) [module, fill=boxgreen!20, below=of d4, minimum width=3.5cm] {差异检测 + 证据产物};
\node (d6) [data, fill=bg, below=of d5, minimum width=3.5cm] {产出业务结论:\\无异动 / 部件差异异动\\图像不匹配异动 / 检测不可用};

\draw [arrow] (d0) -- (d1); \draw [arrow] (d1) -- (d2); \draw [arrow] (d2) -- (d3);
\draw [arrow] (d3) -- (d4); \draw [arrow] (d4) -- (d5); \draw [arrow] (d5) -- (d6);
\end{tikzpicture}
\caption{标定模式（左）和日常检测模式（右）的流程对比}
\label{fig:dual}
\end{figure}


%===========================================================================
\section{完整数据流总结}

\begin{figure}[htbp]
\centering
\footnotesize
\begin{tikzpicture}[node distance=0.5cm]
\node [data] (in) {标准图.jpg + 实时图.jpg + YOLO坐标.txt};
\node [module, fill=boxblue!20, below=of in] (mod1) {references.py: 发现文件、读图\\roi.py: 解析坐标、归一化容错};
\node [data, below=of mod1] (d1) {CaseInput + Roi列表};
\node [module, fill=boxorange!20, below=of d1] (mod2) {alignment.py: ORB→RANSAC→投影几何→Warp→ECC\\ORB不可靠时可选SuperPoint+LightGlue补点→同一套三态诊断};
\node [data, below=of mod2] (d2) {AlignmentResult: H矩阵, aligned\_live, valid\_mask};
\node [module, fill=boxpurple!20, below=of d2] (mod3) {service.py: warpPerspective(live→标准坐标)};
\node [data, below=of mod3] (d3) {aligned\_live \quad (标准坐标系下的实时图)};
\node [module, fill=boxpurple!20, below=of d3] (mod4) {mapping.py: ROI投影→在框比例→有效重叠→分辨率\\→原始NCC+SSIM→三次抗光照复核};
\node [gate, below=of mod4] (g1) {through?};
\node [module, fill=boxred!20, below=0.8cm of g1] (mod5) {compare.py: 抗光照Lab颜色+边缘方向\\+清晰度检查→候选分级};
\node [module, fill=boxgreen!20, below=of mod5] (mod6) {artifacts.py: 5张证据PNG};
\node [data, below=of mod6] (out) {四类结论 + 候选框 + 证据路径};

\node [module, fill=boxred!10, right=2cm of g1] (fail) {检测不可用或图像不匹配异动\\+ 证据产物};

\draw [arrow] (in) -- (mod1);
\draw [arrow] (mod1) -- (d1);
\draw [arrow] (d1) -- (mod2);
\draw [arrow] (mod2) -- (d2);
\draw [arrow] (d2) -- (mod3);
\draw [arrow] (mod3) -- (d3);
\draw [arrow] (d3) -- (mod4);
\draw [arrow] (mod4) -- (g1);
\draw [arrow] (g1) -- node[right] {\footnotesize 通过} (mod5);
\draw [arrow] (g1.east) -- ++(0.3,0) |- node[right] {\footnotesize 不过} (fail);
\draw [arrow] (mod5) -- (mod6);
\draw [arrow] (mod6) -- (out);

% multi-component wrapper annotation
\node [draw, dashed, fit=(in)(mod1)(d1)(mod2)(d2)(mod3)(d3)(mod4)(mod5)(mod6)(out), inner sep=8pt, label={[align=left, font=\footnotesize]right:{multi\_component.py:\\选参考图+逐部件+整图聚合}}] {};
\end{tikzpicture}
\caption{完整数据流（从上到下，每个箭头是一层产出）}
\label{fig:dataflow}
\end{figure}


%===========================================================================
\section{附录：关键指标解释}

\begin{itemize}
\item \textbf{内点率（inlier\_rate）}：RANSAC给出的内点数 ÷ 特征匹配数。越高说明"一致的特征点越多"，配准越可信
\item \textbf{重投影误差（reprojection\_error）}：内点被$H$映射后，跟目标位置的偏差（像素）。越低越好
\item \textbf{空间覆盖度（spatial\_coverage）}：内点在标准图上的分布面积占比。如果所有内点挤在一个角落，就算内点率很高也可能配错
\item \textbf{有效重叠率（valid\_overlap\_ratio）}：warp后有效像素占标准图的比例。黑边太多说明旋转/缩放太大
\item \textbf{ECC相关性（ecc\_correlation）}：微调后两图的全局相关程度。越高说明微调更成功
\item \textbf{有效分辨率尺度（effective\_resolution\_scale）}：$H$的雅可比矩阵最小奇异值——实时图上一个像素相当于标准图上多少个像素。$<1$说明实时图更糙
\item \textbf{比较倍率（comparison\_scale）}：$\min(1.0, S)$，其中 $S$ 是有效分辨率尺度，确保比较时不用上采样

\item \textbf{标定模式 vs 日常检测模式}：标定模式可用开发配置，只产观察数据和证据图，\textbf{不产业务结论}。日常检测模式必须在入口加载带版本号的标定配置，产四类业务结论：高置信无异动、部件差异异动、图像不匹配异动、检测不可用。这是整个系统的"安全闸"——防止用占位阈值冒充生产判定
\end{itemize}

\end{document}
